\documentclass[letterpaper]{article}
\usepackage[submission]{aaai2027}
\usepackage[hyphens]{url}
\usepackage{graphicx}
\urlstyle{rm}
\def\UrlFont{\rm}
\usepackage{natbib}
\usepackage{caption}
\usepackage{booktabs}
\usepackage{amsmath}
\frenchspacing
\setcounter{secnumdepth}{0}

\pdfinfo{
/TemplateVersion (2027.1)
}

\title{Technical Supplement for ``Dual-Expert Keys for Reducing False-Positive Retrieval in Lifelong Vision--Language Model Editing''}
\author{Anonymous Submission}
\affiliations{}

\begin{document}
\maketitle

\section{Scope}
This document lists implementation details that were omitted or compressed in the main paper.
All scientific claims, method equations for Dual-Expert Keys and adversarial visual binding, evaluation protocol sizes, and primary result tables remain in the main text.
Reviewers are not required to consult this supplement.

\section{Editor Configuration}
Primary experiments use the BLIP-2 OPT-2.7B LiveEdit configuration in \texttt{configs/liveedit/blip2-opt-2.7b.yaml}.
Relevant fields not tabulated in the main paper:

\begin{center}
\small
\begin{tabular}{@{}ll@{}}
\toprule
Field & Value \\
\midrule
Retrieval module dimension & $1024$ \\
Cross-attention heads & $8$ \\
Editing query tokens (\texttt{eqe\_n}) & $4$ \\
LoRA rank / scale & $4$ / $5$ \\
LLM mid dimension & $2560$ \\
Edit layer index & $21$ (OPT decoder) \\
Base learning rate & $1\times10^{-4}$ \\
LR cut iteration / rate & $10000$ / $0.1$ \\
Rel / Gen / Loc / soft / hard routing weights & all $1$ \\
\bottomrule
\end{tabular}
\captionof{table}{BLIP-2 LiveEdit YAML settings used for reported runs.}
\end{center}

When patch grounding is enabled, vision logits mix Q-Former and DINOv2 patch scores with fuse weight $\alpha{=}0.5$ (\texttt{patch\_routing\_alpha}).
The DINOv2 encoder is frozen; only the PatchGroundingModule and MoE generators are trained for the DEK/AVBT path.

\section{Probe-Builder Exact Rules}
The conceptual image-swap and hard-negative probes are defined in the main paper.
Here we record the exact pairing rules used by \texttt{tools/build\_grounding\_eval\_json.py}.

\paragraph{Prompt skeleton.}
For a source question string $s$, normalize whitespace, lowercase, and take the first seven tokens.
This skeleton indexes candidates for hard-negative preference.

\paragraph{Entity exclusion.}
\begin{itemize}
\item VLKEB (\texttt{entity\_mode=folder}): the entity id is the first path component of the image relative path.
\item EVQA / EasyEdit (\texttt{entity\_mode=coco}): the entity id is the numeric COCO image id parsed from the filename.
\end{itemize}
Partners must have a different entity id than the source edit.

\paragraph{Ranking.}
Among other-entity rows, assign skeleton score $2$ for an exact skeleton match, $1$ if the candidate appears in the same skeleton bucket list, else $0$.
Hard-negative candidates require score ${>}0$ when such partners exist; otherwise any other-entity row is allowed.
Swap candidates are all other-entity rows, sorted by the same score then by index distance.
The first candidate whose image file exists on disk is kept; hard-negative and swap partners are preferably distinct.

\paragraph{Drop rule.}
A row is dropped if no valid hard-negative or no valid swap partner image exists.

\paragraph{Packaged VLKEB eval overlay stats.}
For the shipped \texttt{data/VLKEB/grounding/grounding\_eval\_short.json} metadata: pairing seed $42$; $3174$ input rows retained as $3174$ outputs; skeleton-match rates $0.696$ (hard-negative) and $0.636$ (swap); zero missing partners after ranking.

\section{FPR Logging Procedure}
The main paper defines false-positive retrieval (FPR) as the fraction of adversarial queries that activate at least one foreign expert.
Operationally, with a built lifelong pool:

\begin{enumerate}
\item Apply sequential edits so each sample contributes one or more pool entries; maintain a map from pool index to sample id.
\item For each image-swap or hard-negative probe, call routing diagnostics (\texttt{probe\_routing}) and record retrieved pool indices.
\item Mark the probe as a false-positive if any retrieved index maps to a sample id other than the query's own sample.
\item Aggregate the indicator over the evaluated adversarial probes (optionally separately for swap vs.\ hard-negative).
\end{enumerate}

In the code supplement this is exposed via \texttt{test\_grounding\_eval.py --log\_routing} and \texttt{tools/log\_grounding\_diagnostics.py}.
Answer-level swap/hard-negative token accuracy is computed separately by the standard evaluation harness and is not used as the FPR definition.

\section{AVBT Auxiliary Terms}
Let $\mathcal{L}_{\mathrm{pos}}$, $\mathcal{L}_{\mathrm{neg}}$, $\lambda$, $\beta$, and hinge $\tau$ be as in the main paper Method section.
On an adversarial image, write $\pi^{\mathrm{adv}}$ for the softmax over pool experts plus the prototype slot, and let
\[
\pi^{\mathrm{wrong}}_{\max}
=
\max_{i\notin\mathcal{O}\cup\{\mathrm{prot}\}}\pi^{\mathrm{adv}}_i
\]
be the largest probability on a non-own expert.
The informal $\mathcal{L}_{\mathrm{aux}}$ used in primary runs (with \texttt{--avbl\_contrastive}) is
\begin{align}
\mathcal{L}_{\mathrm{contrast}}
&=
\mathrm{ReLU}\!\left(\pi^{\mathrm{wrong}}_{\max}-\gamma\right),
\qquad \gamma{=}0.1,
\\
\mathcal{L}_{\mathrm{targeted}}
&=
-\log\!\left(1-\pi^{\mathrm{wrong}}_{\max}\right),
\end{align}
and $\mathcal{L}_{\mathrm{aux}}=\mathcal{L}_{\mathrm{contrast}}+\mathcal{L}_{\mathrm{targeted}}$.
Neighbor-mining and pool-buffer extras are disabled in the primary Dual-Expert Keys row (they appear only as negative ablations in the main paper).

\section{Training Schedule Details}
Reported Dual-Expert Keys checkpoints are obtained by fine-tuning from a lifelong LiveEdit baseline trained on original VLKEB edits (no DEK, no AVBT).
Typical command-line settings for the DEK+AVBT stage:
\begin{itemize}
\item batch size $4$; one training epoch over the probe-augmented VLKEB train JSON;
\item random seed $42$ for the headline run;
\item \texttt{-lkpt} set to the baseline checkpoint path;
\item \texttt{--patch\_grounding --adversarial\_grounding --avbl\_contrastive --dek\_dual\_gate};
\item binding hyperparameters as in the main paper ($\lambda{=}0.3$, $\beta{=}2.5$, $\tau{=}0.3$, text margin $m{=}0$).
\end{itemize}
Model weights, VLKEB/EVQA images, and trained checkpoints are not redistributed with the optional code archive; they must be obtained separately.

\section{Code and Data Artifacts}
Relative paths in the anonymous code and data package:
\begin{itemize}
\item \texttt{configs/liveedit/blip2-opt-2.7b.yaml} --- editor hyperparameters above;
\item \texttt{tools/build\_grounding\_eval\_json.py} --- probe construction;
\item \texttt{data/VLKEB/grounding/} and \texttt{data/easy-edit-mm/grounding/} --- probe-augmented JSON metadata (images excluded);
\item \texttt{train\_vllm\_editor.py} --- training entry point;
\item \texttt{test\_grounding\_eval.py} --- adversarial evaluation and optional routing logs;
\item \texttt{editor/vllm\_editors/liveedit/} --- DEK / AVBT / patch-grounding implementation.
\end{itemize}

\end{document}
